\documentclass[11pt,aps,preprint,nofootinbib,superscriptaddress,tightenlines]{revtex4-2}
\usepackage{amsmath,amssymb,bm,mathtools,graphicx}
\usepackage{xcolor}
\definecolor{darkblue}{RGB}{0,0,160}
\usepackage[colorlinks=true,linkcolor=darkblue,citecolor=darkblue,urlcolor=darkblue]{hyperref}

\def\st{\begin{equation}}
\def\stp{\end{equation}}
\def\Eq#1{eq.~(\ref{#1})}
\def\Sect#1{Sect.~\ref{#1}}
\newcommand{\bbeta}{\bm\beta}
\let\oldparagraph\paragraph
\renewcommand{\paragraph}{\par\bigskip\oldparagraph}

\begin{document}

\title{The ideal step in the energy representation}
\author{D.~Teaney}
\affiliation{Department of Physics and Astronomy, Stony Brook University,
Stony Brook, New York 11794, USA}
\date{\today}

\begin{abstract}
The Hamiltonian is written as the sum of the particle energies,
$H=\sum_ae_a$, each evaluated from the particle's own entropy, momentum and
volume.  The volume is coarse grained with the entropy weights already in the
code, $V_a=s_a/S(\bm x_a)$.  Hamilton's equations then give an ideal step in
which each particle moves with its own velocity and feels only a pressure force.
Viscosity is left out.
\end{abstract}

\maketitle

%=============================================================================
\section{Variables and the volume}
\label{sec:vars}

Each particle $a$ has a position $\bm x_a$, a momentum $\bm m_a$ and an entropy
weight $s_a$.  The weights are constants of the ideal step; a convenient start
is $s_a=1$ for all particles.  The coarse-grained entropy density at the
particle is the kernel sum
\st
S_a\equiv S(\bm x_a)=\sum_bs_b\,W(\bm x_a-\bm x_b),
\stp
and the particle's volume is defined by
\st
V_a=\frac{s_a}{S_a},\qquad\text{so that}\qquad \frac1{V_a}=\frac{S_a}{s_a}
\;\xrightarrow{\,s_a=1\,}\;\sum_bW(\bm x_a-\bm x_b).
\label{Va}
\stp
For unit weights $1/V_a$ is the number density of particles.  The particle's
entropy density $s_a/V_a$ is then the smoothed one, $S_a$.  Its momentum
density is its own, $\bm m_a/V_a$, not smoothed.

When the particles move, the volume changes as
\st
\frac{\dot V_a}{V_a}=-\frac{\dot S_a}{S_a}
=\frac{1}{S_a}\sum_bs_b\,(\dot{\bm x}_b-\dot{\bm x}_a)\cdot\bm D_{ab},
\qquad
\bm D_{ab}\equiv\frac{\partial W(\bm x_a-\bm x_b)}{\partial\bm x_a},
\label{Vdot}
\stp
the SPH estimate of $\nabla\cdot\dot{\bm x}$ at $\bm x_a$.

%=============================================================================
\section{The particle energy}
\label{sec:energy}

The equation of state gives the energy density $h(S,\bm M)$ as a function of the
entropy and momentum densities.  The particle's energy is
\st
e_a=E_a(s_a,\bm m_a,V_a)\equiv V_a\,h\Big(\frac{s_a}{V_a},\frac{\bm m_a}{V_a}\Big),
\label{ea}
\stp
the same particle energy as in the particle representation of
\texttt{changerep.tex}.  Its differential is the first law for the particle,
\st
de_a=T_a\,ds_a+\bm v_a\cdot d\bm m_a-p_a\,dV_a ,
\label{firstlaw}
\stp
where $T_a$ (bare, $T_{\rm rel}/\gamma$), $\bm v_a$ and $p_a$ are the temperature,
velocity and pressure at the particle's densities $(S_a,\bm m_a/V_a)$.

%=============================================================================
\section{The Hamiltonian and Hamilton's equations}
\label{sec:ham}

The Hamiltonian is the sum of the particle energies,
\st
\boxed{\;H(\bm x,\bm m)=\sum_ae_a=\sum_aV_a\,h\Big(\frac{s_a}{V_a},\frac{\bm m_a}{V_a}\Big),
\qquad V_a=\frac{s_a}{S(\bm x_a)}\;}
\label{H}
\stp
It depends on the positions only through the volumes.  Since $s_a=V_aS_a$,
\Eq{H} is the same as
\st
H=\sum_as_a\,\mathcal E\Big(S_a,\frac{\bm m_a}{s_a}\Big),
\stp
with $\mathcal E$ the energy per unit entropy.  The present Hamiltonian
(\texttt{sph2.tex}) is $\sum_as_a\mathcal E(S_a,\bm M_a/S_a)$, with the smoothed
momentum $\bm M_a=\sum_b\bm m_bW_{ab}$ in place of the particle's own
$\bm m_a$; that is the only difference.

Hamilton's equations are
\st
\dot{\bm x}_a=\frac{\partial H}{\partial\bm m_a},\qquad
\dot{\bm m}_a=-\frac{\partial H}{\partial\bm x_a},\qquad \dot s_a=0 .
\stp
By \Eq{firstlaw}, $\partial H/\partial\bm m_a=\bm v_a$: the particle moves with
its own velocity,
\st
\boxed{\;\dot{\bm x}_a=\bm v_a\;}
\label{xdot}
\stp
For the force, $\partial H/\partial\bm x_a=-\sum_cp_c\,\partial V_c/\partial\bm x_a$ with
$\partial V_c/\partial\bm x_a=-(V_c/S_c)\,\partial S_c/\partial\bm x_a$, and
\st
\frac{\partial S_c}{\partial\bm x_a}=\delta_{ca}\sum_bs_b\bm D_{ab}+s_a\bm D_{ac}.
\stp
Collecting terms,
\st
\boxed{\;\dot{\bm m}_a=-\sum_bs_as_b\left(\frac{p_a}{S_a^2}+\frac{p_b}{S_b^2}\right)\bm D_{ab}\;}
\label{mdot}
\stp
This is the standard symmetric SPH pressure force, with $S$ in the role of the
mass density and $s_a$ in the role of the mass.  Compared with \texttt{sph2.tex},
the terms in $\Delta\mathcal M_b(\bm x_a)=\bm m_b/s_b-\bm M_a/S_a$ are gone,
because the momentum is not smoothed.

\paragraph{Continuum limit.}  With $\sum_bV_bf_b\bm D_{ab}\to\nabla f$,
\Eq{mdot} becomes $\dot{\bm m}_a\to-V_a\nabla p$.  This is the momentum equation
for a parcel moving with $\bm v$: relative to such a parcel the momentum flux
$T^{ij}=M^iv^j+p\,\delta^{ij}$ reduces to $p\,\delta^{ij}$.

\paragraph{The particle energies.}  From \Eq{firstlaw} at fixed $s_a$,
\st
\dot e_a=\bm v_a\cdot\dot{\bm m}_a-p_a\dot V_a ,
\label{edot}
\stp
with $\dot V_a$ from \Eq{Vdot}.  The particle energies change by the work done on
the particle; their sum does not.

%=============================================================================
\section{What is conserved}
\label{sec:cons}

\begin{enumerate}
\item \emph{Energy.}  $H=\sum_ae_a$ is conserved by Hamilton's equations, and it is
  the same sum of particle energies \Eq{ea} that the viscous step of
  \texttt{changerep.tex} works with: the smoothed and the particle energies are
  one and the same.
\item \emph{Momentum.}  $\bm D_{ba}=-\bm D_{ab}$, so the pair forces in
  \Eq{mdot} cancel and $\sum_a\bm m_a$ is conserved.
\item \emph{Entropy.}  Each $s_a$ is constant, so $\sum_as_a$ is conserved, and
  each particle's entropy $S_a(e_a,\bm m_a)$ at volume $V_a$ equals $s_a$.
\end{enumerate}

%=============================================================================
\section{What it does not change: the static rigidity}
\label{sec:rigidity}

At rest, $\bm m_a=0$, the new and the present Hamiltonians are the same function of
the positions,
\st
H\big|_{\bm m=0}=\sum_as_a\,\mathcal E(S_a,0),
\stp
because the two differ only in which momentum enters $\mathcal E$.  The static phonon
spectrum of \texttt{rigidity.tex} (Sect.~``The ideal Hamiltonian's phonons''),
and with it the static shear modulus $G$ and $c_t$, are therefore unchanged: the
positions still enter only through the kernel sum $S_a=\sum_bs_bW_{ab}$, which is
what breaks the relabeling symmetry.  The energy representation is \emph{not} a
fix for the rigidity.  (Volumes that partition space, $\sum_aV_a=A$, would be; that
is a different choice of $V_a$ in \Eq{Va}, and \Eq{H}--\Eq{mdot} would carry over
with $\partial V_c/\partial\bm x_a$ from the partition.)

What does change is the dynamics: the particles move with their own velocities
rather than a doubly smoothed one, and the force has no $\Delta\mathcal M$ terms.
Whether that changes the measured transverse oscillation in $C_T$ (its frequency
or damping) is an open question for the runs, not something this note predicts.
The gain that is certain is consistency with the viscous step, \Sect{sec:visc}.

%=============================================================================
\section{The deterministic viscous step}
\label{sec:visc}

With $H=\sum_ae_a$, the viscous step is done on the particles' own variables; there
is no change of representation.  The derivative of $H$ with respect to a weight at
fixed $\bm x$, $\bm m$ is
\st
T^{\rm eff}_a\equiv\frac{\partial H}{\partial s_a}
=T_a-\frac{p_a}{S_a}+\sum_bs_b\frac{p_b}{S_b^2}W_{ab}
=T_a+O(h^2),
\label{Teff}
\stp
using $\mathcal E-\bm v\cdot\bm m/s=T-p/S$; the last two terms are the smoothed minus
the local $p/S$.  With
\st
\bbeta_a=\frac{\bm v_a}{T^{\rm eff}_a},\qquad
\Delta\bm m_a=-\Delta t\,\frac{\partial K}{\partial\bbeta_a},\qquad
\Delta s_a=-\bbeta_a\cdot\Delta\bm m_a ,
\label{viscstep}
\stp
$\Delta H=\sum_a(T^{\rm eff}_a\Delta s_a+\bm v_a\cdot\Delta\bm m_a)=0$ to first order
in $\Delta t$, and $\sum_a\Delta s_a=2K\Delta t\ge0$, as in \texttt{sph2.tex}.  $K$ is
the bond dissipative function of \texttt{changerep.tex}, which already takes a
pointwise $\bbeta_a$.

Compared with \texttt{changerep.tex}: the forward map $s_a\to e_a$ is no longer a
map, since $e_a$ is already the energy in $H$; the gradient correction $g_a$
(its Appendix~A) was the difference between $\sum_ae_a$ and the smoothed $H$, and is
now zero.  The exact back map at fixed $e_a$, $V_a$ does not carry over as it is: a
change of $s_a$ changes $S_b$, hence $V_b$, at every neighbour, so fixing every
$e_a$ at fixed $V_a$ does not fix $H$ exactly.  The step \Eq{viscstep} conserves $H$
to first order in $\Delta t$; the drift is measured (plan, step 6).

%=============================================================================
\section{Implementation plan}
\label{sec:plan}

The Hamiltonian is chosen by a new field \texttt{params.hamiltonian} $\in$
\{\texttt{:smoothed}, \texttt{:energy}\} (environment \texttt{SPH\_HAMILTONIAN}),
default \texttt{:smoothed}, so that every existing run and test is unchanged.

\paragraph{Unchanged.}  The \texttt{Particle} struct (\texttt{x}, \texttt{s},
\texttt{m}, \texttt{S\_smooth} and the EOS fields are enough); kernels, cell and
neighbour lists; the EOS solvers \texttt{eos\_newton}, \texttt{solve\_eos!},
\texttt{particle\_energy}, \texttt{eos\_from\_energy}; the integrators
\texttt{leapfrog\_step!} and \texttt{kdk\_step!} (their structure; they call the
functions below); the bond dissipative function and force; the set-up of initial
conditions ($s_a=S(\bm x_a)\,{\rm vol}$, or $s_a=1$); the recorders.  The stochastic
step is left alone.

\paragraph{Changed (all in \texttt{celllist.jl}, dispatched on \texttt{params.hamiltonian}).}
\begin{enumerate}
\item \texttt{smooth\_to\_particles!}: sum $S_a$ only, and set the momentum density to
  the particle's own, \texttt{M\_smooth} $=\bm m_aS_a/s_a=\bm m_a/V_a$.  Then
  \texttt{solve\_eos!} returns, unchanged, the particle's own $T_a$, $\bm v_a$, $p_a$
  and $e_a/s_a$ (\texttt{H\_over\_S}).  One neighbour sum fewer.
\item \texttt{update\_velocities!}: $\dot{\bm x}_a=\bm v_a$, local, no neighbour sum;
  \texttt{Teff} from \Eq{Teff}, one neighbour sum (needed only by the viscous step and
  by the bond metric).
\item \texttt{hamiltonian\_force!}: the pair bracket becomes
  $-s_as_b(p_a/S_a^2+p_b/S_b^2)$, \Eq{mdot}.
\item Viscous step: \texttt{dissipative\_step\_changerep!} with $\bbeta_a$ of
  \Eq{viscstep} read from the particle (no \texttt{pointwise\_state} copy), and
  $\Delta s_a=-\bbeta_a\cdot\Delta\bm m_a$.  \texttt{dissipative\_step!} calls it, not
  the gradient-corrected step, when \texttt{params.hamiltonian == :energy}.
\item Diagnostics: the total energy is $\sum_ae_a=\sum_as_a\,$\texttt{H\_over\_S}$_a$,
  the same number as $H$.
\end{enumerate}

\paragraph{Steps and checks.}
\begin{enumerate}
\item Implement 1--3 behind the flag; \texttt{test/testregress\_step.jl} must stay
  bit-identical for \texttt{:smoothed}.
\item A test \texttt{test/testenergystep.jl}: $\partial H/\partial\bm m_a$ and
  $\partial H/\partial\bm x_a$ against finite differences of $H$ on a jittered
  lattice; $\sum_a\dot{\bm m}_a=0$; $\dot s_a=0$.
\item Ideal runs: energy drift of $H$ versus $\Delta t$ (second order for the
  leapfrog), a sound wave (speed and damping) against the present step
  (\texttt{examples/testsound.jl}), and the time step: the unsmoothed $\bm v_a$ may
  need a smaller CFL.
\item Static check: the phonon spectrum of \Sect{sec:rigidity} must agree with the
  present one to roundoff (it is the same $H$ at $\bm m=0$).
\item Implement 4; check $\sum_a\Delta s_a=2K\Delta t$ and $\Delta H=O(\Delta t^2)$
  per step; a shear wave against $e^{-\nu k^2t}$ (\texttt{examples/testshearwave.jl}).
\item Energy drift over long coupled deterministic runs, against the present step with
  the gradient-corrected map.
\end{enumerate}

\paragraph{Documentation.}  Each changed function gets a comment that states the
equation it implements and cites this note by equation; \texttt{params.hamiltonian}
is described where \texttt{SimParams} is defined; the outcome of each check above
is added to this note as it is done.

\paragraph{Status (2026-11-09).}  Items 1--3 of the changes and item 4 (the viscous step,
\texttt{dissipative\_step\_energy!}) are in \texttt{src/sph.jl} and
\texttt{src/celllist.jl}.  On a jittered lattice with a moving fluid, $\partial H/\partial\bm
m_a$, $\partial H/\partial\bm x_a$ and $\partial H/\partial s_a$ agree with $\dot{\bm x}_a$,
$-\dot{\bm m}_a$ and $T^{\rm eff}_a$ to $10^{-6}$ (finite differences); $\sum_a\dot{\bm
m}_a=7\times10^{-16}$; $H=\sum_a$\texttt{particle\_energy} exactly.  The deterministic
shear wave is \Sect{sec:shear}.

%=============================================================================
\section{The deterministic shear wave}
\label{sec:shear}

\texttt{examples/shearwave\_decay.jl} (viscous and ideal) and
\texttt{examples/shearwave\_scan.jl} (ideal, kernels and $h/a$).  A uniform fluid with
\st
M_y(x,0)=A\sin kx,\qquad k=\frac{2\pi}{L},
\stp
evolved with the full deterministic step.  A fluid keeps the amplitude with $\eta=0$
and decays as $e^{-\nu k^2t}$, $\nu=(\eta/s)/T_{\rm rel}$, with $\eta>0$.  An elastic
solid with transverse speed $c_t$ instead gives, with $\eta=0$,
\st
\frac{A(t)}{A(0)}=\cos(c_tkt).
\stp
At $h/a=2$ (C2, $kh=0.31$):
\begin{enumerate}
\item The two Hamiltonians agree: amplitudes to $10^{-3}$ with and without viscosity,
  $|\Delta H/H|\sim10^{-11}$ for the ideal step.  The viscous step drifts by
  $4\times10^{-8}$ (\texttt{:energy}) and $8\times10^{-9}$ (\texttt{:smoothed}).
\item For $A\le0.005$ (displacements $\le0.4a$) the ideal wave is an undamped cosine with
  $c_t=0.0340$ (\texttt{:energy}) and $0.0335$ (\texttt{:smoothed}): the static value
  $0.034$ of \texttt{rigidity.tex}.  For $A=0.02$ the particles slip by $2.4a$ and the
  response is not linear.
\item With $\eta/s=0.02$ the amplitude falls below $e^{-\nu k^2t}$ and turns negative: a
  damped transverse oscillation.
\end{enumerate}

\paragraph{Kernels and $h/a$.}  \texttt{examples/shearwave\_scan.jl}: ideal step only,
$A=0.002$, a periodic strip $L_x=1$, $L_y\simeq6h$ (2D, the wave depends on $x$ only;
560 to 3400 particles), $t\le30$.  $c_t$ is a least-squares fit of $\cos(c_tkt)$, or
$\cosh(|c_t|kt)$ for an unstable wave (written $c_t<0$), over the first quarter period.
Dynamic $c_t$ (\texttt{:energy} / \texttt{:smoothed}) against the static phonon values of
\texttt{rigidity.tex} ($qh\simeq0.26$; here $kh=0.31$):
\begin{center}
\begin{tabular}{llccccc}
\hline
kernel & & $h/a=2$ & $2.5$ & $3$ & $4$ & $5$ \\ \hline
C2 & dynamic & $0.0340/0.0335$ & $0.0308/0.0304$ & $0.0214/0.0211$ & $-0.005^{\,a}$ & $-0.0027/-0.0027$ \\
   & static  & $0.035$ & $0.031$ & $0.022$ & $-0.011$ & $-0.0034$ \\
C4 & dynamic & grows$^{\,b}$ & $-0.0047/-0.0047$ & $-0.0018/-0.0018$ & $0.0022/0.0022$ & $0.0008/0.0008$ \\
   & static  & $-0.0054$ & unstable & $-0.0029$ & $0.0029$ & $0.0008$ \\
C6 & dynamic & $0.0170/0.0169$ & $0.0065/0.0064$ & $0.0027/0.0026$ & $0.0004/0.0004$ & $0.0002/0.0002$ \\
   & static  & $0.017$ & $0.0066$ & $0.0026$ & $-0.0003$ & $0.0002$ \\
\hline
\end{tabular}
\end{center}
$^a$ grows, then not linear (poor fit).  $^b$ $A$ rises to $1.03$ by $t\simeq10$, then the
displacement reaches $0.9a$ and the response is not linear.

\begin{enumerate}
\item The dynamic and static $c_t$ agree, signs included, wherever the wave stays linear: the
  transverse sound of the runs is the elastic stiffness of $H$, and the time stepping adds
  nothing.
\item The two Hamiltonians agree to $1$--$2\%$, as they must ($H$ is the same function of
  the positions at $\bm m=0$).  The energy representation by itself does not change the
  rigidity.
\item $|c_t|$ falls with $h/a$ and with the kernel order, oscillating in sign (C2, C4); C6
  is below $10^{-3}$ from $h/a=4$ on.  $\Delta H/H\le10^{-13}$ in every linear run.
\end{enumerate}
(The strips at $h/a=3$ and $5$ have two cell rows in $y$; they first blew up through a
double counting in the cell list for \texttt{ncy}${}=2$, fixed in \texttt{neighbor\_range}.)

%=============================================================================
\section{Proposal: volumes that add up to the area}
\label{sec:norm}

\paragraph{Why the kernel volume gives a modulus.}  Since $H$ depends on the positions only
through the volumes, to second order in their changes
\st
\delta H=-\sum_ap_a\,\delta V_a+\sum_a\frac{K_a}{2V_a}\,\delta V_a^2 ,
\label{dH}
\stp
with $K$ the bulk modulus.  For $V_a=s_a/S_a$ a homogeneous shear $\epsilon$ of a
lattice changes every volume by the same amount at second order, $\delta V_a=c\,\epsilon^2$,
so $\sum_a\delta V_a\neq0$, and the first term gives a modulus
\st
G\sim p\,c(h/a,\text{kernel}).
\stp
This is the $\epsilon$ of \texttt{rigidity.tex}, proportional to the pressure.

\paragraph{The condition.}  If the volumes always add up to the area,
\st
\sum_aV_a=A ,
\label{sumV}
\stp
then at uniform pressure $\sum_ap\,\delta V_a=p\sum_a\delta V_a=0$, and the second term of
\Eq{dH} is of fourth order in $\epsilon$ on a lattice: no modulus.  A partition of space
(Voronoi cells, or the partition of unity of meshless finite-mass methods) satisfies
\Eq{sumV}; so does the simplest choice, a global normalization of the kernel volume,
\st
V_a=\frac{A}{Z}\,\tilde V_a,\qquad \tilde V_a=\frac{s_a}{S_a},\qquad Z=\sum_c\tilde V_c .
\label{Vnorm}
\stp
The particle's entropy density is then $s_a/V_a=(Z/A)\,S_a$.

\paragraph{Hamilton's equations.}  With \Eq{Vnorm} in \Eq{H},
$\partial H/\partial\bm x_a=-\sum_cp_c\,\partial V_c/\partial\bm x_a$ and
\st
\frac{\partial V_c}{\partial\bm x_a}=\frac AZ\Big(\frac{\partial\tilde V_c}{\partial\bm x_a}
-\frac{\tilde V_c}{Z}\sum_d\frac{\partial\tilde V_d}{\partial\bm x_a}\Big).
\stp
The second term subtracts the volume-averaged pressure
\st
\bar p=\frac1A\sum_cp_cV_c ,
\stp
and the force is \Eq{mdot} with $p\to p-\bar p$:
\st
\dot{\bm m}_a=-\frac AZ\sum_bs_as_b\left(\frac{p_a-\bar p}{S_a^2}+\frac{p_b-\bar p}{S_b^2}\right)\bm D_{ab}.
\label{mdotnorm}
\stp
The same substitution in \Eq{Teff} gives the derivative with respect to the weights,
\st
T^{\rm eff}_a=T_a-\frac AZ\left[\frac{p_a-\bar p}{S_a}-\sum_bs_b\,\frac{p_b-\bar p}{S_b^2}\,W_{ab}\right].
\stp
Still exact: $H$ is conserved by Hamilton's equations, and the pair terms of
\Eq{mdotnorm} are odd under $a\leftrightarrow b$, so momentum is conserved.  The continuum
limit is unchanged, $\nabla(p-\bar p)=\nabla p$.  The cost is one global sum per step.
In SPH, subtracting a background pressure is a known remedy for the tensile
instability; here it follows from the Hamiltonian, and energy conservation is kept.

\paragraph{Expected, and tests.}
\begin{enumerate}
\item Perfect lattice: $c_t$ falls from $0.034$ to near zero.  Test: the shear-wave scan
  with \Eq{Vnorm}, against \Sect{sec:shear}.
\item Thermal state: what is left is driven by the pressure fluctuations $\delta p$
  rather than by $p$, with random sign; no coherent long-wavelength modulus is expected.
  Test: the stochastic runs, \Sect{sec:stoch}.
\item Disordered arrangements: an affine shear changes the individual $\delta V_a$ at first
  order, with random sign, which leaves a modulus $\propto K\,(\text{disorder})^2$ from the
  second term of \Eq{dH}: the caveat (a) of \texttt{rigidity.tex} for Voronoi volumes.  A
  local partition does not remove it either.  Test: the static phonon spectrum on jittered
  lattices (\texttt{ideal\_lattice\_modes.jl}, \texttt{JITTER}).
\end{enumerate}
Implementation: an option \texttt{volume = :normalized} of the \texttt{:energy}
Hamiltonian; one reduction for $\bar p$ and $Z$, and $p\to p-\bar p$, $S\to(Z/A)S$ where
the particle state is evaluated.

%=============================================================================
\section{The stochastic step with \texttt{:energy}}
\label{sec:stoch}

Not yet tested.  The stochastic step (\texttt{stochrep.tex}) works at fixed positions,
volumes $V_a=s_a/S_a$ and particle energies $e_a$, which are exactly the variables of
\Eq{H}; on a frozen lattice it is the same step for both Hamiltonians.  Two differences:
\begin{enumerate}
\item \emph{No gradient correction.}  The gradient energy $g_a$ of \texttt{changerep.tex}
  (Appendix~A) absorbs the difference between $\sum_ae_a$ and the smoothed $H$, which is
  now zero.  With \texttt{:energy} the step must use the plain back map
  $s_a'=S_a(e_a,\bm m_a')$ (\texttt{gradient = false}).
\item \emph{The back map and the volumes.}  The back map holds $e_a$ and $V_a$ fixed, but a
  change of $s_a$ changes $S_b$ and so $V_b$ at every neighbour.  To first order
  \st
  \Delta H\simeq\sum_a\big(T^{\rm eff}_a-T_a\big)\Delta s_a ,
  \stp
  with $T^{\rm eff}_a-T_a$ the smoothed minus the local $p/S$, of order $h^2$.  In the
  deterministic step this is $4\times10^{-8}$ (\Sect{sec:shear}); with noise the $\Delta s_a$
  are larger and fluctuate, so its size has to be measured.
\end{enumerate}
Tests, in order:
\begin{enumerate}
\item \texttt{stochrep\_step!} defaults to \texttt{gradient = false} when
  \texttt{params.hamiltonian == :energy}, and the drivers keep that default.
\item One step on a thermal state: $\Delta H$ per stochastic step, \texttt{:energy} against
  \texttt{:smoothed} (with its gradient-corrected map).
\item Coupled runs, \texttt{stochrep\_hamiltonian.jl}: $L=0.6$, $h/a=2$, $\eta/s=0.2$,
  $S_0=10^6$, to $t\approx50$, 2--4 realizations per Hamiltonian: $\Delta H/H(t)$,
  KE/eq, the lattice disorder.  Local.
\item At NERSC, after 1--3: equilibrium $C_T$, $C_L$ with \texttt{:energy}, as campaign
  \texttt{second}; the node-hour estimate first.
\end{enumerate}
\paragraph{Status.}  Test 1 is done; \texttt{:smoothed} stays bit-identical
(\texttt{test/testregress\_step.jl}).  Test 2 (\texttt{examples/energystep\_stoch\_onestep.jl},
$L=0.6$, $h/a=2$, equilibrium start, 20 half steps $\Delta t/2=0.0043$ at fixed positions,
3 seeds):
\begin{center}
\begin{tabular}{lcc}
\hline
 & \texttt{:smoothed} (gradient map) & \texttt{:energy} \\ \hline
$\Delta H/H$ per half step, rms & $3.4\times10^{-5}$ & $1.8\times10^{-7}$ \\
$\Delta H/H$ per half step, mean & $\pm\,7\times10^{-6}$ (seed to seed) & $-1.8\times10^{-7}$ (every seed) \\
\hline
\end{tabular}
\end{center}
With \texttt{:energy}, $\Delta H=\Delta\sum_ae_a$ exactly, and the change is $200$ times smaller
than with \texttt{:smoothed}, but systematic: a cooling of about $4\times10^{-5}$ per unit
time.  It is correlated ($-0.8$) with $\sum_a(T^{\rm eff}_a-T_a)\Delta s_a$, but of the
opposite sign and twice the size; not understood yet.

With \Sect{sec:norm}, the same tests with \texttt{volume = :normalized} decide whether the
transverse oscillation of the equilibrium $C_T$ is the pressure term of \Eq{dH}.

\end{document}
